\documentclass[11pt]{article}

\usepackage[a4paper,margin=24mm]{geometry}
\usepackage{amsmath,amssymb,amsthm,mathtools}
\usepackage{subfiles}
\usepackage{xcolor}
\usepackage{microtype}
\usepackage[pdfusetitle,colorlinks=true,linkcolor=blue,citecolor=blue,
urlcolor=blue]{hyperref}
\usepackage{interval}
\intervalconfig{soft open fences}
\usepackage{derivative}
\usepackage{mleftright}
\usepackage{tikz}
\usepackage{graphicx}
\usetikzlibrary{arrows.meta,positioning,calc}

\newtheorem{theorem}{Theorem}[section]
\newtheorem{lemma}[theorem]{Lemma}
\newtheorem{proposition}[theorem]{Proposition}
\newtheorem{corollary}[theorem]{Corollary}
\theoremstyle{remark}
\newtheorem{remark}[theorem]{Remark}

\newcommand{\E}{\mathbb E}
\newcommand{\Var}{\operatorname{Var}}
\newcommand{\tr}{\operatorname{tr}}
\newcommand{\I}{\mathbb I}
\newcommand{\N}{\mathcal N}
\newcommand{\calI}{\mathcal I}
\newcommand{\norm}[1]{\left\lVert #1\right\rVert}

\title{A Central Limit Theorem for High-Temperature SYK Model: A case study on mathematical research using AI and Lean}
\author{Yoshinori Kamijima\thanks{Information Networking for Innovation and Design, Toyo University. E-mail: kamijima@iniad.org}
\and Shuta Nakajima\thanks{Department of Mathematics, Keio University. E-mail: njima@keio.jp}
\and Ryoichiro Noda\thanks{Faculty of Science and Engineering, Waseda University.
E-mail: r\textunderscore noda@aoni.waseda.jp}
\and K\^{o}hei Sasaya\thanks{Graduate School of Mathematical Sciences, The University of Tokyo. E-mail: sasaya@g.ecc.u-tokyo.ac.jp}
}
\date{}

\begin{document}

\maketitle

\begin{abstract}
The Sachdev--Ye--Kitaev (SYK) model is a quantum many-body model of Majorana
fermions with all-to-all random interactions. It has attracted wide interest in condensed matter and high-energy physics as a solvable model of non-Fermi liquids and holography. Fix an even integer $q\geq4$, and let $Z_N(\beta)$ be the
partition function of the Gaussian $q$-body SYK model.  We prove that for all sufficiently small inverse temperatures
$\beta > 0$, there exists $\rho=\rho(\beta)>0$ such that
\[
N^{q/2 -1}(\log Z_N(\beta)-\E[\log Z_N(\beta)])
  \Longrightarrow \N(0,\rho^2).
\]
In addition to the mathematical results, we discuss how the proofs were developed and refined through interaction between human researchers and AI, and we describe the subsequent formalization of the main arguments in Lean 4.
This work thus serves as a case study of a research workflow that integrates human insight, AI-assisted exploration, iterative proof refinement, and formal verification.
\end{abstract}

\noindent\emph{Keywords:} Sachdev--Ye--Kitaev model,
free energy, central limit theorem, Gaussian chaos.

\medskip
\noindent\emph{Mathematics Subject Classification (2020):}
60F05, 60B20, 82B44.



\subfile{sections/introduction}
\subfile{sections/gaussian-chaos}
\subfile{sections/higher-chaoses}
\subfile{sections/main-results}

\appendix
\section*{Lean formalization.}
We have formalized the proofs of our main results in Lean 4, conditional
on Proposition~\ref{prop:annealed-free-energy}. Proposition~\ref{prop:annealed-free-energy} is
treated as an explicit assumption and has not
been formalized. 
The Lean code is available at
\url{https://github.com/TokyoProbabilityGroup/SYK_CLT}.

\section*{Acknowledgements}
  The first author was supported by JSPS KAKENHI Grant Number JP26K17003.
 The second author was supported by JSPS KAKENHI Grant Numbers JP24K16937 and JP25K00911. 
 The third author was supported by JSPS KAKENHI Grant Number JP26KJ0336. 
  The fourth author was supported by JSPS KAKENHI Grant Numbers JP24KJ0022 and JP26K17012.
\begin{thebibliography}{99}

\bibitem{AGK2025}
E.~R. Anschuetz, D.~Gamarnik, and B.~T. Kiani,
\newblock Bounds on the ground state energy of quantum $p$-spin Hamiltonians,
\newblock \emph{Commun. Math. Phys.} 406 (2025), article 232.
\newblock \href{https://doi.org/10.1007/s00220-025-05412-4}
{doi:10.1007/s00220-025-05412-4}.



\bibitem{Billingsley1999}
P.~Billingsley,
\newblock \emph{Convergence of Probability Measures},
\newblock 2nd ed., Wiley, New York, 1999.
\newblock \href{https://doi.org/10.1002/9780470316962}
{doi:10.1002/9780470316962}.

\bibitem{FTW2018}
R.~Feng, G.~Tian, and D.~Wei,
\newblock Spectrum of SYK model II: Central limit theorem,
\newblock \emph{Random Matrices Theory Appl.} 10 (2021), no.~4, 2150037.
\newblock \href{https://doi.org/10.1142/S2010326321500374}
{doi:10.1142/S2010326321500374}.

\bibitem{GPSZ2026}
D.~Gamarnik, F.~Pernice, A.~Schmidhuber, and A.~Zlokapa,
\newblock The free energy limit of the SYK model at high temperature,
\newblock \href{https://arxiv.org/abs/2605.02768}
{arXiv:2605.02768v1}, 2026.

\bibitem{GGTW2025}
B.~Georgiev, J.~G\'omez-Serrano, T.~Tao, and A.~Z.~Wagner,
\newblock Mathematical exploration and discovery at scale,
\newblock \href{https://arxiv.org/abs/2511.02864}
{arXiv:2511.02864}, 2025.

\bibitem{He2026}
Y.~He,
\newblock Free energy and spectral edge of the SYK model,
\newblock \href{https://arxiv.org/abs/2607.18998}
{arXiv:2607.18998v3}, 2026.


\bibitem{HiaiPetz2014}
F.~Hiai and D.~Petz,
\newblock \emph{Introduction to Matrix Analysis and Applications},
\newblock Universitext, Springer (2014).
\newblock \href{https://doi.org/10.1007/978-3-319-04150-6}
{doi:10.1007/978-3-319-04150-6}.

\bibitem{Hoeffding1963}
W.~Hoeffding,
\newblock Probability inequalities for sums of bounded random variables,
\newblock \emph{J. Amer. Statist. Assoc.} 58 (1963), no.~301, 13--30.
\newblock \href{https://doi.org/10.1080/01621459.1963.10500830}
{doi:10.1080/01621459.1963.10500830}.


\bibitem{Kitaev2015}
A.~Kitaev,
\newblock A simple model of quantum holography (part 2),
\newblock Talk at the KITP program \emph{Entanglement in Strongly-Correlated
Quantum Matter}, May 27, 2015.
\newblock \url{https://online.kitp.ucsb.edu/online/entangled15/kitaev2/}.

\bibitem{Macdonald1995}
I.~ G.~ Macdonald,
\newblock \emph{Symmetric Functions and Hall Polynomials},
\newblock 2nd ed.,
\newblock Oxford Mathematical Monographs,
\newblock Oxford University Press, Oxford, 1995.

\bibitem{MaldacenaStanford2016}
J.~Maldacena and D.~Stanford,
\newblock Remarks on the Sachdev--Ye--Kitaev model,
\newblock \emph{Phys. Rev. D} 94 (2016), 106002.
\newblock \href{https://doi.org/10.1103/PhysRevD.94.106002}
{doi:10.1103/PhysRevD.94.106002}.

\bibitem{OpenAI}
OpenAI.
\newblock Sharing AI progress in mathematics.
\newblock 2026.
\url{https://openai.com/index/sharing-ai-progress-in-mathematics/}
(accessed October 21, 2026).


\bibitem{Rockafellar1970}
R.~T. Rockafellar,
\newblock \emph{Convex Analysis},
\newblock Princeton University Press, 1970.

\bibitem{SachdevYe1993}
S.~Sachdev and J.~Ye,
\newblock Gapless spin-fluid ground state in a random quantum Heisenberg
magnet,
\newblock \emph{Phys. Rev. Lett.} 70 (1993), 3339--3342.
\newblock \href{https://doi.org/10.1103/PhysRevLett.70.3339}
{doi:10.1103/PhysRevLett.70.3339}.

\bibitem{S2023}
D.~W.~Stroock,
\newblock \emph{Gaussian Measures in Finite and Infinite Dimensions},
\newblock Universitext, Springer (2023).
\newblock \href{https://doi.org/10.1007/978-3-031-23122-3}
{doi:10.1007/978-3-031-23122-3}.



\bibitem{Zlokapa2026}
A.~Zlokapa,
\newblock A rigorous quasipolynomial-time classical algorithm for SYK thermal
expectations,
\newblock \href{https://arxiv.org/abs/2604.21089}
{arXiv:2604.21089v1}, 2026.
\end{thebibliography}

\end{document}
